\documentclass[aspectratio=169,11pt]{beamer}
% Shared Beamer style for practice contest solution walkthroughs.
\usepackage[utf8]{inputenc}
\usepackage[T1]{fontenc}
\usepackage{lmodern}
\usepackage{amsmath,amssymb}
\usepackage{xcolor}
\usepackage{hyperref}
\usepackage{listings}

\definecolor{CPNavy}{HTML}{172033}
\definecolor{CPBlue}{HTML}{2563EB}
\definecolor{CPGreen}{HTML}{16A34A}
\definecolor{CPOrange}{HTML}{F97316}
\definecolor{CPGray}{HTML}{64748B}
\definecolor{CPLight}{HTML}{F8FAFC}
\definecolor{CPCodeBg}{HTML}{F1F5F9}
\definecolor{CPCodeRule}{HTML}{CBD5E1}

\usetheme{default}
\usefonttheme{professionalfonts}
\setbeamertemplate{navigation symbols}{}
\setbeamersize{text margin left=0.65cm,text margin right=0.65cm}
\setbeamertemplate{itemize item}{\color{CPBlue}$\blacktriangleright$}
\setbeamertemplate{itemize subitem}{\color{CPGray}$\triangleright$}

\setbeamercolor{normal text}{fg=CPNavy,bg=white}
\setbeamercolor{frametitle}{fg=white,bg=CPNavy}
\setbeamercolor{title}{fg=white,bg=CPNavy}
\setbeamercolor{subtitle}{fg=CPLight,bg=CPNavy}
\hypersetup{colorlinks=true,urlcolor=CPBlue}

\setbeamertemplate{frametitle}{%
  \nointerlineskip%
  \begin{beamercolorbox}[wd=\paperwidth,ht=0.88cm,dp=0.22cm,leftskip=0.55cm,rightskip=0.35cm]{frametitle}
    \usebeamerfont{frametitle}\insertframetitle\hfill\small\insertframenumber/\inserttotalframenumber
  \end{beamercolorbox}%
}

\setbeamertemplate{footline}{%
  \leavevmode%
  \hbox{%
    \begin{beamercolorbox}[wd=.58\paperwidth,ht=2.4ex,dp=1ex,leftskip=.5cm]{author in head/foot}%
      \color{CPGray}\insertshorttitle
    \end{beamercolorbox}%
    \begin{beamercolorbox}[wd=.42\paperwidth,ht=2.4ex,dp=1ex,rightskip=.5cm plus1fil]{date in head/foot}%
      \hfill\color{CPGray}\insertshortauthor
    \end{beamercolorbox}%
  }%
}

\setbeamertemplate{title page}{%
  \vspace*{1.25cm}
  \begin{beamercolorbox}[wd=\paperwidth,sep=0.65cm,leftskip=0.15cm,rightskip=0.15cm]{title}
    {\color{white}\Huge\bfseries\inserttitle\par}
    \vspace{0.35cm}
    {\color{CPLight}\Large\insertsubtitle\par}
    \vspace{0.6cm}
    {\color{CPLight}\large\insertauthor\par}
  \end{beamercolorbox}
  \vfill
}

\newcommand{\sectiontag}[1]{\textbf{\color{CPBlue}#1}}
\newenvironment{tightitemize}{\begin{itemize}\small\setlength{\itemsep}{0.18cm}}{\end{itemize}}
\lstdefinestyle{cpcode}{
  language=Python,
  basicstyle=\ttfamily\tiny,
  keywordstyle=\color{CPBlue}\bfseries,
  commentstyle=\color{CPGray}\itshape,
  stringstyle=\color{CPGreen},
  backgroundcolor=\color{CPCodeBg},
  frame=single,
  framerule=0.25pt,
  rulecolor=\color{CPCodeRule},
  showstringspaces=false,
  columns=fullflexible,
  keepspaces=true,
  breaklines=true,
  breakatwhitespace=false,
  tabsize=4,
  xleftmargin=0.08cm,
  xrightmargin=0.08cm,
  aboveskip=0.08cm,
  belowskip=0.08cm
}


\title[snapperhard]{Snapper Chain (Hard)}
\subtitle{Day 1 solution walkthrough}
\author[Solution Deck]{Competitive Programming Bootcamp}
\date{}

\begin{document}

\begin{frame}[plain]
  \titlepage
\end{frame}

% BEGIN GENERATED STATEMENT INTERPRETATION
\begin{frame}{Interpret the Word Problem}
\small
\sectiontag{Statement excerpts:}
\begin{tightitemize}
  \item \textit{``Snapper devices and chained them together''}
  \item \textit{``Will the light be on or off''}
\end{tightitemize}

\vspace{0.18cm}
\sectiontag{Programming interpretation:}
\begin{tightitemize}
  \item The chain state follows the low bits of a binary counter.
  \item The light is on exactly when all N snappers are on and powered.
  \item That condition is equivalent to the lowest N bits of K all being 1.
\end{tightitemize}
\end{frame}
% END GENERATED STATEMENT INTERPRETATION







\begin{frame}{Problem Model}
\small
\sectiontag{Textbook connection:} CP4 2.2 h / 2.2.3 Bitmask

\vspace{0.25cm}
\sectiontag{Core technique:} Bitmask pattern recognition

\vspace{0.25cm}
\begin{tightitemize}
  \item The chain behaves like a binary counter after each clap.
  \item The light is on only when all N snappers are on.
  \item That means the lowest N bits of K must all be 1.
\end{tightitemize}
\end{frame}

\begin{frame}{How to Recognize the Approach}
\begin{tightitemize}
  \item Simulating K claps is unnecessary and may be too slow.
  \item The repeated states reveal powers of two.
  \item A mask with N low bits set tests the condition in constant time.
\end{tightitemize}
\end{frame}

\begin{frame}{Algorithm}
\begin{tightitemize}
  \item For each case, build mask = (1 << N) - 1.
  \item Compute K \& mask to isolate the lowest N bits.
  \item If the isolated bits equal mask, print ON.
  \item Otherwise print OFF.
\end{tightitemize}
\end{frame}

\begin{frame}{Walkthrough of the Key State}
\begin{tightitemize}
  \item For N = 3, mask is binary 111.
  \item K \& 111 discards all higher bits of K.
  \item Only values ending in 111 have all three snappers on.
\end{tightitemize}
\end{frame}

\begin{frame}{Why the Algorithm Is Correct}
\begin{tightitemize}
  \item The snapper chain advances through the same states as binary counting.
  \item A snapper is powered only if every lower snapper is on, so the final light depends on all N low bits.
  \item The bitmask check is exactly the condition that all those bits are 1.
\end{tightitemize}
\end{frame}

\begin{frame}{Complexity and Implementation Notes}
\small
\sectiontag{Complexity:} O(T) time and O(1) extra memory for T test cases.

\vspace{0.3cm}
\begin{tightitemize}
  \item Parentheses around 1 << N make the mask expression clear.
  \item This is a good example of replacing simulation with a state-pattern observation.
\end{tightitemize}
\end{frame}



% BEGIN GENERATED CODE WALKTHROUGH
\begin{frame}[fragile]{Code: Read Test Cases}
\scriptsize
\sectiontag{Source lines:} \texttt{snapperhard.py:39--44}

\vspace{0.08cm}
\begin{columns}[T,totalwidth=\textwidth]
\begin{column}{0.58\textwidth}
\begin{lstlisting}[style=cpcode]
import sys

data = sys.stdin.read().strip().splitlines()

for idx, case in enumerate(data[1:]):
    n, k = list(map(int, case.split()))
\end{lstlisting}
\end{column}
\begin{column}{0.38\textwidth}
\sectiontag{What this chunk does}
\begin{tightitemize}
  \item The first input line is the number of cases.
  \item Each case supplies N, the chain length, and K, the number of snaps.
  \item enumerate gives the 1-based case labels used in the output.
\end{tightitemize}
\end{column}
\end{columns}
\end{frame}

\begin{frame}[fragile]{Code: Test the Low Bits}
\scriptsize
\sectiontag{Source lines:} \texttt{snapperhard.py:45--49}

\vspace{0.08cm}
\begin{columns}[T,totalwidth=\textwidth]
\begin{column}{0.58\textwidth}
\begin{lstlisting}[style=cpcode]
mask = (1 << n) - 1  # check if all lower bits below n are turned on
if (k & mask) == mask:  # all n low bits are 1
    print(f"Case #{idx+1}: ON")
else:
    print(f"Case #{idx+1}: OFF")
\end{lstlisting}
\end{column}
\begin{column}{0.38\textwidth}
\sectiontag{What this chunk does}
\begin{tightitemize}
  \item The mask contains N ones in binary.
  \item k \& mask isolates exactly the N snappers that matter.
  \item If those bits equal the mask, every snapper in the chain is ON.
\end{tightitemize}
\end{column}
\end{columns}
\end{frame}
% END GENERATED CODE WALKTHROUGH

\begin{frame}{Source}
\small
\sectiontag{Solution file:} \texttt{practice\_contests/day\_1/snapperhard.py}

\vspace{0.3cm}
\sectiontag{Problem:} \url{https://open.kattis.com/problems/snapperhard}

\vspace{0.45cm}
Use this deck with the source code open beside it. The slides explain the reasoning; the code shows the exact input handling and output formatting.
\end{frame}

\end{document}
